\documentclass[11pt,reqno]{amsart}



\usepackage[utf8]{inputenc}
\usepackage{mathtools,amsthm,amssymb}
\usepackage[allowunsupported,grade=19]{cmgraded}
\usepackage{microtype}
\usepackage[a4paper,top=25mm,bottom=26mm,left=29mm,right=29mm,headsep=7mm]{geometry}
\usepackage[dvipsnames]{xcolor}
\usepackage[
  breaklinks=true,
  pagebackref=true,
  hypertexnames=false
]{hyperref}
\hypersetup{
  colorlinks=true,
  pdfpagemode=UseNone,
  citecolor=ForestGreen,
  linkcolor=MidnightBlue,
  urlcolor=RoyalBlue,
  pdfstartview=FitH,
}


% ===== 自行调整的参数 =====
\linespread{1.05}           % 行距倍率；1 为 amsart 原始行距
\normalsize                  % 使行距倍率立即生效
\linespacing=\baselineskip   % 同步章节、定理使用的间距基准
\setlength{\parindent}{1em} % 段首缩进
\setlength{\parskip}{3pt}   % 段间距

\makeatletter

% ===== 首页标题区 =====
\def\@settitle{%
  \begin{center}%
    \uppercasenonmath\@title
    \Large\@title
  \end{center}}

\def\@setauthors{%
  \vskip21pt plus6pt minus3pt
  \begin{center}%
    {\tiny BY\par}\medskip
    {\normalsize\scshape\authors\par}\medskip
    \def\address##1##2{\par ##2\par}%
    \def\email##1##2{\par e-mail:~##2\par}%
    {\small\itshape\addresses\par}%
  \end{center}}

\def\enddoc@text{} % 地址已经排在首页

\renewenvironment{abstract}{%
  \global\setbox\abstractbox=\vtop\bgroup
  \vskip10pt
  \normalfont\tiny\centering\MakeUppercase{\abstractname}\par
  \Small
  \list{}{%
    \leftmargin3pc
    \rightmargin3pc
    \labelwidth0pt
    \listparindent\normalparindent
    \itemindent0pt
    \parsep0pt plus1pt
    \let\fullwidthdisplay\relax
  }%
  \item[]
}{%
  \vspace*{-10pt}
  \endlist\egroup
}

% ===== 摘要下方的 AI 使用声明 =====
\newbox\aiusagebox
\newenvironment{aiusage}{%
  \global\setbox\aiusagebox=\vtop\bgroup
  \normalfont\fontsize{8}{10.5}\selectfont
  \list{}{%
    \leftmargin3pc
    \rightmargin3pc
    \topsep0pt
    \partopsep0pt
    \parsep0pt
    \itemsep0pt
    \listparindent0pt
    \itemindent0pt
  }%
  \item[]%
  \noindent{\color{black!25}\rule{\linewidth}{0.3pt}}\par
  \nobreak\vspace{4pt}%
  {\normalfont\tiny\centering\MakeUppercase{AI usage}\par}%
  \nobreak\vspace{2pt}%
  \noindent\ignorespaces
}{%
  \par\endlist\egroup
}
\let\original@setabstracta\@setabstracta
\def\@setabstracta{%
  \original@setabstracta
  \ifvoid\aiusagebox\else
    \vskip10pt
    \box\aiusagebox
    \prevdepth\z@
  \fi
}

% ===== 章节层级 =====
\let\@secnumfont\@empty

\def\section{\@startsection{section}{1}{0pt}%
  {20pt plus4pt minus4pt}
  {.5\linespacing}
  {\normalfont\bfseries}}

\def\subsection{\@startsection{subsection}{2}{0pt}%
  {.5\linespacing plus.7\linespacing}
  {-.5em}
  {\normalfont\scshape}}

\def\subsubsection{\@startsection{subsubsection}{3}{0pt}%
  {.5\linespacing plus.7\linespacing}
  {-.5em}
  {\normalfont\itshape}}

% 章节标题后的第一段不缩进
\let\@afterindenttrue\@afterindentfalse
\@afterindentfalse

% ===== 定理、定义、备注 =====
\newtheoremstyle{plain}
  {.5\linespacing}{.5\linespacing}
  {\itshape}{1em}{\scshape}{.}{.5em}
  {\thmname{#1}\thmnumber{ #2}\thmnote{\normalfont{} (#3)}}

\theoremstyle{plain}
\newtheorem{theorem}{Theorem}[section]
\newtheorem{lemma}[theorem]{Lemma}
\newtheorem{proposition}[theorem]{Proposition}
\newtheorem{corollary}[theorem]{Corollary}

\newtheoremstyle{definition}
  {.5\linespacing}{.5\linespacing}
  {\upshape}{1em}{\scshape}{.}{.5em}
  {\thmname{#1}\thmnumber{ #2}\thmnote{\normalfont{} (#3)}}

\theoremstyle{definition}
\newtheorem{definition}[theorem]{Definition}

\newtheoremstyle{remark}
  {.5\linespacing}{.5\linespacing}
  {\upshape}{1em}{\scshape}{.}{.5em}
  {\thmname{#1}\thmnumber{ #2}\thmnote{\normalfont{} (#3)}}

\theoremstyle{remark}
\newtheorem{remark}[theorem]{Remark}

% ===== 证明与结束标记 =====
\renewenvironment{proof}[1][\proofname]{%
  \par
  \pushQED{\qed}%
  \normalfont
  \topsep6pt plus6pt\relax
  \trivlist
  \item[\hskip\labelsep\slshape #1\@addpunct{.}]%
  \ignorespaces
}{%
  \popQED
  \endtrivlist
  \@endpefalse
}

\DeclareRobustCommand{\qed}{%
  \leavevmode\unskip\penalty9999\hbox{}\nobreak
  \ifmmode
    \makebox[0pt][l]{\qquad\hbox{\qedsymbol}}%
  \else
    \qquad\hbox{\qedsymbol}%
  \fi
}
\renewcommand{\qedsymbol}{\rule[.47pt]{4pt}{6pt}}

\makeatother
 


\title[Speed monotonicity of BRW on GW trees]{More rootward bias, less speed:\\
 biased walks on BGW trees  in the B\"{o}ttcher case}
\author{Heng Ma}
\address[Heng Ma]{Faculty of Data and Decision Sciences, Technion,  Israel}
\email{hengmamath(at)gmail(dot)com}
\urladdr{\url{https://hengmamath.github.io}}
\date{\today}


\setlength{\emergencystretch}{1.8em}
\numberwithin{equation}{section}
\allowdisplaybreaks[1]
 

\newcommand{\E}{\mathbb E}
\newcommand{\PP}{\mathbb P}
\newcommand{\Law}{\operatorname{Law}}
\newcommand{\Cov}{\operatorname{Cov}}
\newcommand{\dd}{\,\mathrm d}
\newcommand{\Dplus}{D^{+}}
\newcommand{\cJ}{\mathcal J}
\newcommand{\cL}{\mathcal L}
\newcommand{\cF}{\mathcal F}
\newcommand{\res}{\mathcal R}

\begin{document}


\begin{abstract}
Let $v(\lambda)$ be the speed of the $\lambda$-biased random walk on a
Bienaym\'e--Galton--Watson tree with finite offspring mean $m$.
If every vertex has at least two children, we prove that $v$ is strictly
decreasing throughout the transient regime $[0,m)$.
We give two proofs based on a common two-child reweighting estimate:
one uses a correction that vanishes at recursive fixed points, and the
other compares escape laws by stochastic order and freezes the descendants
to isolate the root response.
\end{abstract}

\begin{aiusage}
The initial proof was found by OpenAI's
GPT-6 Astra through extensive interaction with the author. During the conversation, the model frequently lost in short-term rewards:
it proved local lemmas that did not lead to the
main result, checked special cases numerically, or carried out long
computations without real progress. It also tended to keep pushing one
approach instead of stepping back. In these cases, the author either chose among directions proposed by the
model or asked it to try a different approach.
 The author has independently checked all arguments  and takes full
responsibility for the results.
\end{aiusage}


\maketitle

% 首页的剩余空白留在页底，避免拉大声明与目录之间的间距。
\begingroup
  \raggedbottom
  \vspace*{-20pt}% 压缩 amsart 标题区与目录之间的默认留白。
  \setcounter{tocdepth}{1}% 两条证明共享的主章节列在首页；子节由正文导航。
  \tableofcontents
  \clearpage
\endgroup

\section{Introduction}\label{sec:introduction}

Let $\mathcal{T}$ be a Galton--Watson tree, and let $\xi$ have its offspring
law $(p_k)_{k\ge0}$. We assume throughout that
\begin{equation}\label{eq:assumptions}
 \PP(\xi\ge2)=1,\qquad m:=\E\xi<\infty.
\end{equation}
At a vertex with $k$ children, the $\lambda$-biased walk jumps to its parent
with probability $\lambda/(\lambda+k)$ and to each child with probability
$1/(\lambda+k)$. Thus increasing $\lambda$ increases the bias towards the root.
Let $W_n$ be its position at time $n$. The convention for transitions from
the root does not affect the speed; see \cite[Lemma~2.1]{SWX}. For
$0<\lambda<m$, write
\begin{equation}\label{eq:speed-def}
 v(\lambda)=\lim_{n\to\infty}\frac{|W_n|}{n}.
\end{equation}
The limit exists and is a deterministic positive constant in this range
\cite[Section~1]{Aidekon}. At $\lambda=0$, the walk moves away from the root
at every step, so we set $v(0)=1$.

Lyons, Pemantle and Peres asked whether $v(\lambda)$ is nonincreasing in $\lambda$
for Galton--Watson trees without leaves \cite[Question~2.1]{LPP}.
The increase in the probability of moving to the parent at each vertex
does not settle the question: changing $\lambda$ also changes the parts of
the random tree visited by the walk. Song, Wang and Xiang proved strict
decrease on an interval $[0,a_{m_1}]$ with $a_{m_1}>1$ when the minimal
offspring number $m_1$ is at least two \cite[Theorem~1.2]{SWX}. For the
uniform offspring law on $\{2,3\}$, Mandarapu and Kunkunuru subsequently
verified strict decrease up to $\lambda=351/200$ by a computer-assisted
argument \cite[Theorem~1]{MK}. The result below covers the entire
positive-speed range for every finite-mean law satisfying
\eqref{eq:assumptions}.

For a real-valued function $f$, write its upper right Dini derivative as
\begin{equation}\label{eq:dini-def}
 \Dplus f(\lambda)
 :=\limsup_{\varepsilon\downarrow0}
       \frac{f(\lambda+\varepsilon)-f(\lambda)}{\varepsilon}.
\end{equation}

\begin{theorem}\label{thm:main}
Under \eqref{eq:assumptions}, for every $1<\lambda<m$ there exists
$\delta_\lambda>0$ such that
\begin{equation}\label{eq:main-dini}
 \Dplus v(\lambda)\le-\delta_\lambda<0.
\end{equation}
Consequently, $v(\lambda)$ is strictly decreasing in $\lambda$ on $[0,m)$.
\end{theorem}

\begin{corollary}\label{cor:derivative}
If, in addition, $\E e^{\theta\xi}<\infty$ for some $\theta>0$, then
$v(\lambda)$ is differentiable in $\lambda$ on $(0,m)$ and
\begin{equation}\label{eq:negative-derivative}
 v'(\lambda)<0,\qquad 0<\lambda<m.
\end{equation}
\end{corollary}

Both proofs establish \eqref{eq:main-dini} on $(1,m)$, then use the
overlapping low-parameter interval of \cite{SWX} to obtain the global
conclusion. The assumption of at least two children enters a single
reweighting estimate; only the finite mean is needed in the fixed-point
comparison, so the offspring number may be unbounded. The arguments do not
give the corresponding estimate when vertices with one child are allowed.
The additional exponential-moment assumption in Corollary~\ref{cor:derivative} is
used only to invoke the differentiability theorem of Bowditch and
Tokushige \cite[Theorem~1.1]{BT}.

\subsection{Proof strategy}

Adjoin an artificial parent to the root, and let $\beta_\gamma$ be the
probability that a walk started at the root never visits this parent. Its
annealed law $\pi_\gamma=\Law(\beta_\gamma)$ satisfies the recursive equation
\[
 X\ \overset{\mathrm d}=\ f_\gamma\Bigl(\sum_{i=1}^{\xi}X_i\Bigr),
 \qquad f_\gamma(s)=\frac{s}{\gamma+s},
\]
where the $X_i$ are independent with law $\pi_\gamma$ and independent of
$\xi$. These laws decrease stochastically as $\gamma$ increases. A
symmetric form of A\"id\'ekon's speed formula, proved in
Proposition~\ref{prop:symmetric}, is
\begin{equation}\label{eq:roadmap-speed}
 v(\gamma)
 =\frac{\E[XY/(\gamma-1+X+Y-XY)]}
        {\E[(X+Y-XY)/(\gamma-1+X+Y-XY)]},
 \qquad 1<\gamma<m,
\end{equation}
where $X,Y$ are independent with law $\pi_\gamma$. Write
$N_\gamma(\rho)$ and $Z_\gamma(\rho)$ for the numerator and denominator
when both variables instead have law $\rho$.

Fix $1<\lambda<m$. The sign of $v(\gamma)-v(\lambda)$ is that of
\[N_\gamma(\pi_\gamma)-v(\lambda)Z_\gamma(\pi_\gamma).\]

Both the kernels and the
law change with $\gamma$, and stochastic ordering of the laws does not
by itself determine the sign. For a smooth function $u$, the recursive
residual
\begin{equation}\label{eq:roadmap-residual}
 \res_u(\gamma,\rho)
 =\E_\rho u(X)-\E_{\xi,\rho}
   u\Bigl(f_\gamma\bigl(\textstyle\sum_{i=1}^{\xi}X_i\bigr)\Bigr)
\end{equation}
vanishes at every actual fixed point $\rho=\pi_\gamma$. We introduce the
corrected functional
\begin{equation}
  \cF(\gamma,\rho)=N_\gamma(\rho)-v(\lambda)Z_\gamma(\rho)-2\res_u(\gamma,\rho).
\end{equation}
At $\rho=\pi_\gamma$, the residual vanishes, giving
\[
 \cF(\gamma,\pi_\gamma)
 =Z_\gamma(\pi_\gamma)\{v(\gamma)-v(\lambda)\}.
\]
We will choose $u$ to obtain, for $\gamma>\lambda$ sufficiently close to
$\lambda$, the two comparisons
\[
 \cF(\gamma,\pi_\gamma)
 \le\cF(\gamma,\pi_\lambda)
 <\cF(\lambda,\pi_\lambda)=0.
\] 

 
We choose a single $u$ so that $\cF$ is locally increasing in the input
law under stochastic order and, with that law fixed at $\pi_\lambda$,
$\partial_\gamma\cF(\lambda,\pi_\lambda)<0$. Both $u$ and $v(\lambda)$
are kept fixed in these variations.
To describe the first requirement, let $U\sim\pi_\lambda$
and choose $V$ on the same probability space with $U\le V\le1$ almost surely. Dependence within the pair $(U,V)$ is allowed; and the marginal law of $V$ is arbitrary.
  Let
$\rho_s=\Law(U+s(V-U))$ for $0\le s\le1$. 

For any smooth $u$, with the parameter fixed at $\lambda$, the right
derivative of $s\mapsto\cF(\lambda,\rho_s)$ at $s=0$ has the form
\[
 \left.\frac{\mathrm d}{\mathrm ds}\cF(\lambda,\rho_s)\right|_{s=0+}
 =2\E\bigl[(V-U)G_{\lambda,\pi_\lambda}(U)\bigr].
\]
The coefficient $G_{\lambda,\pi_\lambda}$ depends on the choice of $u$;
the derivative formula is proved in Lemma~\ref{lem:transport}.
Since $V-U\ge0$, the first requirement is uniform positivity of this
coefficient.

Put $c=\lambda-1$ and $D_\lambda(x,y)=c+x+y-xy$, and let $Y\sim\pi_\lambda$.
The uncorrected comparison $N_\lambda-v(\lambda)Z_\lambda$ contributes a term
$\E[Y^2/D_\lambda(x,Y)^2]$ to this coefficient. The following choice
cancels that term and also gives the required negative parameter derivative:
\begin{equation}\label{eq:roadmap-corrector}
 u'(x)=\E\frac{Y}{D_\lambda(x,Y)^2}.
\end{equation}
Proposition~\ref{prop:cancel-law} verifies that the recursive residual
correction cancels the quadratic term, leaving
\[
 G_{\lambda,\pi_\lambda}(x)
 =c\,\E\frac{(1+v(\lambda))Y-v(\lambda)}{D_\lambda(x,Y)^2},
 \qquad 0\le x\le1.
\]
Its positivity is equivalent to
$\widehat\mu_1(x)>v(\lambda)/(1+v(\lambda))$, where
$\widehat\mu_1(x)=\frac{\E[Y/D_\lambda(x,Y)^2]}{\E[1/D_\lambda(x,Y)^2]}$
is the normalized weighted mean. Relative to the weighting at $x=0$,
the density is proportional to $\{(c+y)/D_\lambda(x,y)\}^2$, which
increases in $y$. Hence $\widehat\mu_1(x)\ge\widehat\mu_1(0)$ for
every $x\in[0,1]$ (Lemma~\ref{lem:positive-kernel}), and it suffices to
bound $\widehat\mu_1(0)$. Write $\mu=\E\beta_\lambda$. The hypothesis $\xi\ge2$ enters
Lemma~\ref{lem:tilt}, which proves
\begin{equation}\label{eq:roadmap-tilt}
 \mu<2\widehat\mu_1,
 \qquad
 \widehat\mu_1=\frac{\E[\beta_\lambda/(c+\beta_\lambda)^2]}
         {\E[1/(c+\beta_\lambda)^2]}.
\end{equation}
Here $\widehat\mu_1$ is exactly the weighted mean at $x=0$: the reweighting
can lower the original mean $\mu$, but not by a factor of two.
The symmetric speed formula~\eqref{eq:roadmap-speed}, together with a
decreasing-weight covariance bound, gives
\[
 \frac{v(\lambda)}{1+v(\lambda)}\le\frac\mu2<\widehat\mu_1.
\]
Thus the remaining kernel is uniformly positive, meeting the first
requirement. For the same choice of $u$, the second cancellation in
Proposition~\ref{prop:cancel-parameter} uses this positive kernel to give
$\partial_\gamma\cF(\lambda,\pi_\lambda)<0$ with the law held fixed.
The two signs needed for the comparison therefore follow from one correction.

Finally, for $\gamma>\lambda$ close to $\lambda$, take
$U_\gamma=\beta_\gamma(\mathcal{T})$ and $V_\lambda=\beta_\lambda(\mathcal{T})$ on the same
random tree. Their marginal laws are $\pi_\gamma$ and $\pi_\lambda$,
respectively, and $U_\gamma\le V_\lambda$ almost surely. Interpolate
between their laws using this ordered coupling. Continuity keeps the entire
interpolation in the region of positive law variation. Changing the law
and then the parameter gives a finite comparison with a strictly negative
first-order term, proving \eqref{eq:main-dini} without differentiating
$\beta_\lambda$ on the infinite tree. The sections below establish the
symmetric representation, the reweighting bound, the two cancellation
identities, and the final comparison in that order.

\subsection{Proof strategy 2}\label{subsec:proof-strategy-2}

Both proofs convert the branching recursion into an inverse-square
moment constraint and a positive variation kernel. They then use finite
comparisons of nearby ordered laws and differentiate an explicit upper
bound. The first corrects the functional by a recursive residual; the
second uses the original speed ratio and freezes the descendants.

The second proof works directly with the symmetric speed functional
$\mathcal V_\gamma(\rho)=N_\gamma(\rho)/Z_\gamma(\rho)$. Fix
$1<\lambda<m$, and write $\pi=\pi_\lambda$, $v=v(\lambda)$,
$c=\lambda-1$ and $D=D_\lambda$. Let $X,Y$ be independent with law
$\pi$. We compare the actual escape law at
$\lambda+\varepsilon$ with the law obtained by changing the parameter
at the root while keeping all child escape probabilities at their
values for $\lambda$.

For $0<\varepsilon<m-\lambda$, let $S_\gamma$ be the sum of the child
escape probabilities at parameter $\gamma$. The same-tree coupling gives
$S_{\lambda+\varepsilon}\le S_\lambda$, whereas the recursion at
$\lambda$ gives $S_\lambda=\lambda\beta_\lambda/(1-\beta_\lambda)$.
Consequently,
\[
 \beta_{\lambda+\varepsilon}
 \le\frac{S_\lambda}{\lambda+\varepsilon+S_\lambda}
 =T_{\lambda,\varepsilon}(\beta_\lambda),
 \qquad
 T_{\lambda,\varepsilon}(x)
 =\frac{\lambda x}{\lambda+\varepsilon(1-x)}.
\]
Thus $\pi_{\lambda+\varepsilon}\le_{\mathrm{st}}
(T_{\lambda,\varepsilon})_\#\pi$, where the pushforward denotes the
law of $T_{\lambda,\varepsilon}(X)$. To turn this order into a speed
bound, we need the speed functional to increase with the input law
near $\pi$, and its value along the root perturbation to decrease.

Both signs come from a stronger consequence of the common reweighting
estimate. For $X\sim\pi$, let $\widehat\pi$ have density proportional
to $(c+x)^{-2}$ relative to $\pi$, and write
$\mu=\E X$ and $\widehat\mu_j=\E_{\widehat\pi}X^j$. Put
$S=S_\lambda$, and let $\mathbb Q$ be the probability measure on the
root environment obtained by weighting its law with $(c+S)^{-2}$.
The two-child estimate gives
\[
 \E_{\mathbb Q}S\ge2\widehat\mu_1
 \quad\Longrightarrow\quad
 \lambda(\widehat\mu_1-\widehat\mu_2)
 \ge2\widehat\mu_1(1-2\widehat\mu_1+\widehat\mu_2).
\]
The bound $\xi\ge2$, independence under the original law and conditional
exchangeability under the symmetric tilt give the factor two; this is the shared
branching input of the two proofs. The methods separate at the next
step. The first uses this constraint to obtain
$\mu<2\widehat\mu_1$ (Lemma~\ref{lem:tilt}); the second extracts the
stronger inequality
\[
 \Gamma:=\lambda(2\widehat\mu_1-\mu)
       -(1-\mu)(\widehat\mu_1-\widehat\mu_2)>0.
\]
For $\widehat\mu_1\ge1/2$ this inequality is immediate. In the remaining
case, Proposition~\ref{prop:root-Gamma} expresses $\mu$ through the
tilted moments and checks a concave quadratic at the endpoints of the
feasible interval for $\widehat\mu_2$.

The role of $\Gamma$ is most transparent through the function
\[
 \widetilde h_v(y)
 =(1+v)y-v-\frac{1-v}{2\lambda}y(1-y).
\]
The covariance bound $v\le\mu/(2-\mu)$ gives
$\E_{\widehat\pi}\widetilde h_v(X)
\ge\Gamma/[\lambda(2-\mu)]>0$.
Moreover, $\widetilde h_v$ is increasing. The law weighted by
$D(x,Y)^{-2}$, relative to $\widehat\pi$, has density proportional
to $\{(c+y)/D(x,y)\}^2$, which is increasing in $y$.
With $(\cJ g)(x)=\E[g(Y)/D(x,Y)^2]$, we therefore have
\[
 (\cJ\widetilde h_v)(x)
 \ge\frac{\Gamma}{\lambda^3(2-\mu)}>0,
 \qquad 0\le x\le1.
\]

The common variational object is now visible. Set
$h_v(y)=(1+v)y-v$. Both proofs start with the same kernel
\[
 \Psi=\cJ(y^2+c h_v)
\]
for variation of the original speed functional. In the first proof,
the recursive residual cancels $\cJ(y^2)$, leaving
\[
 G_{\lambda,\pi}(x)
 =\Psi(x)-(\cJ(y^2))(x)
 =c(\cJ h_v)(x)>0.
\]
Here $\mu<2\widehat\mu_1$, the bound
$v/(1+v)\le\mu/2$, and the same increasing comparison of weights
give the required positivity. This is the cancellation in
Proposition~\ref{prop:cancel-law}.

The second proof keeps $\Psi$ and controls it through
$\cJ\widetilde h_v>0$. The decomposition
\[
 y^2+c h_v(y)
 =c\widetilde h_v(y)+y^2+
       \frac{c(1-v)}{2\lambda}y(1-y)
\]
shows that $\Psi$ is uniformly positive. This positivity persists in
a joint neighborhood of the parameter $\lambda$ and the law $\pi$.
A quantile interpolation between ordered laws stays in this
neighborhood and makes the derivative of $\mathcal V_\gamma$
nonnegative. This proves the local
stochastic-order comparison of Lemma~\ref{lem:root-A}, called Lemma A.

For the root response, define
\[
 H_\lambda(\varepsilon)
 =\mathcal V_{\lambda+\varepsilon}
       ((T_{\lambda,\varepsilon})_\#\pi).
\]
The velocity of the root transport at $\varepsilon=0$ is
$-x(1-x)/\lambda$. Its contribution combines with the explicit
parameter derivative of the speed ratio. Writing $N=N_\lambda(\pi)$
and $Z=Z_\lambda(\pi)$, the identity $N-vZ=0$ at the base law
cancels the remaining term, giving
\[
 H_\lambda'(0)
 =-\frac{2}{Z_\lambda(\pi)}
       \E[X(1-X)(\cJ\widetilde h_v)(X)]<0.
\]
This is the root-response bound of Lemma~\ref{lem:root-B}, called
Lemma B. The same positive weighted average thus supplies both signs.

Finally, continuity puts $\pi_{\lambda+\varepsilon}$ and
$(T_{\lambda,\varepsilon})_\#\pi$ in the neighborhood of Lemma A for
sufficiently small $\varepsilon>0$. Applying that lemma at the
perturbed parameter yields
\[
 v(\lambda+\varepsilon)
 =\mathcal V_{\lambda+\varepsilon}(\pi_{\lambda+\varepsilon})
 \le H_\lambda(\varepsilon)
 =v(\lambda)+H_\lambda'(0)\varepsilon+o(\varepsilon).
\]
Hence $\Dplus v(\lambda)\le H_\lambda'(0)<0$. Both proofs handle the
changing escape law by finite comparisons, with reference laws
$\pi_\lambda$ and $(T_{\lambda,\varepsilon})_\#\pi_\lambda$,
respectively, and differentiate only the resulting explicit upper bounds.

\section{Escape probabilities and a symmetric speed formula}\label{sec:reduction}

We now justify the representation and the order properties used in the
introduction. The symmetric representation will allow us to compare speeds
through two independent escape probabilities; the order and continuity
will later connect the actual fixed points by a short interpolation.

Adjoin an artificial parent to the root. Let $\beta_\lambda(\mathcal{T})$ denote the
quenched probability of never visiting that parent, starting from the root.
For $0<\lambda<m$, $0<\beta_\lambda<1$ almost surely. With
$\pi_\lambda=\Law(\beta_\lambda)$, the recursive distributional equation is
\begin{equation}\label{eq:rde}
 X\ \overset{\mathrm d}=\ f_\lambda(S),\qquad
 f_\lambda(s)=\frac{s}{\lambda+s},\qquad
 S=\sum_{i=1}^{\xi}X_i,
\end{equation}
where $X_i$ are independent with law $\pi_\lambda$ and independent of $\xi$.
The RDE follows by a first-step decomposition, including repeated returns
from children to the root.

For $\gamma>1$ and a probability measure $\rho$ on $[0,1]$, define
\begin{align}
 q(x,y)&=1-(1-x)(1-y)=x+y-xy,
 & D_\gamma(x,y)&=\gamma-1+q(x,y),\label{eq:kernels}\\
 N_\gamma(\rho)&=\E_\rho\frac{XY}{D_\gamma(X,Y)},
 & Z_\gamma(\rho)&=\E_\rho\frac{q(X,Y)}{D_\gamma(X,Y)},\label{eq:NZ}
\end{align}
Here $X,Y$ are independent with law $\rho$; $q(x,y)$ is the probability
of at least one success in two independent trials with success
probabilities $x,y$, and $D_\gamma$ is the common denominator.
The quantities $N_\gamma$ and $Z_\gamma$ will be the numerator and
normalizing factor in the speed formula.
For $\rho\ne\delta_0$, the denominator is positive, and we write
\begin{equation}\label{eq:speed-functional}
 \mathcal V_\gamma(\rho)=\frac{N_\gamma(\rho)}{Z_\gamma(\rho)}.
\end{equation}

\begin{proposition}[Symmetric representation]\label{prop:symmetric}
For $1<\lambda<m$,
\begin{equation}\label{eq:symmetric-speed}
 v(\lambda)=\frac{N_\lambda(\pi_\lambda)}{Z_\lambda(\pi_\lambda)}
 =\frac{\E[XY/D_\lambda(X,Y)]}
       {\E[q(X,Y)/D_\lambda(X,Y)]}.
\end{equation}
\end{proposition}

\begin{proof}
A\"id\'ekon's formula \cite[Theorem~1.1]{Aidekon} reads
\begin{equation}\label{eq:aidekon}
 v(\lambda)=
 \frac{\E[(\xi-\lambda)X_0/(\lambda-1+X_0+S)]}
      {\E[(\xi+\lambda)X_0/(\lambda-1+X_0+S)]},
\end{equation}
where $X_0,X_1,\ldots$ are independent copies of $\beta_\lambda$ and
$S=\sum_{i=1}^{\xi}X_i$. Conditional exchangeability gives
\begin{equation}\label{eq:exchange-speed}
 \E\frac{\xi X_0}{\lambda-1+X_0+S}
 =\E\frac{S}{\lambda-1+X_0+S}.
\end{equation}
Set $X=X_0$ and $Y=S/(\lambda+S)$. By \eqref{eq:rde}, $X,Y$ are independent
with law $\pi_\lambda$, and
\begin{equation}\label{eq:denom-change}
 \lambda-1+X+S=\frac{D_\lambda(X,Y)}{1-Y}.
\end{equation}
The numerator in \eqref{eq:aidekon} is consequently
\[
 \lambda\E\frac{Y-X+XY}{D_\lambda(X,Y)}
 =\lambda N_\lambda(\pi_\lambda),
\]
since the $Y-X$ term vanishes by symmetry. The denominator is
$\lambda Z_\lambda(\pi_\lambda)$.
\end{proof}

\begin{lemma}[Order and continuity of the fixed-point curve]\label{lem:continuity}
For $0<\gamma_1<\gamma_2<m$, escape probabilities can be coupled on the same tree
so that $\beta_{\gamma_2}\le\beta_{\gamma_1}$. Moreover,
$\lambda\mapsto\pi_\lambda$ is continuous in the $1$-Wasserstein metric.
In particular, $v$ is continuous on $(1,m)$.
\end{lemma}

\begin{proof}
For a finite-depth tree, impose escape value $1$ at the terminal generation
and iterate the map in \eqref{eq:rde} towards the root. The map is increasing
in every child value and decreasing in the parameter. Induction gives the
claimed order at every depth. The finite-depth probabilities decrease to
the probability of escaping the parent forever, giving the same order in
the limit.

The annealed escape probability $\E\beta_\lambda$ is continuous by
\cite[Proposition~2.9]{BT}. In the leafless case the proof uses
\cite[Lemma~4.7]{BT}, without an exponential-moment assumption. For ordered
random variables, the coupling just constructed gives
\begin{equation}\label{eq:W1}
 W_1(\pi_{\gamma_1},\pi_{\gamma_2})
 =\E\beta_{\gamma_1}-\E\beta_{\gamma_2}.
\end{equation}
Thus the fixed-point laws are $W_1$-continuous. The kernels in
\eqref{eq:NZ} are bounded and continuous locally in $\gamma>1$, and the
normalizing denominator is positive at a nonzero fixed point. This proves
the stated continuity of $v$.
\end{proof}

\section{The common two-child reweighting estimate}\label{sec:tilt}

Both proofs begin with the same constraint on the first two moments
under inverse-square reweighting. The first proof extracts a bound for
the weighted mean and uses a correction of the speed comparison. The
second extracts a stronger scalar inequality and compares the original
speed functional directly. We establish the common constraint once.

Fix $\lambda>1$, put $c=\lambda-1$, and let $X\in(0,1)$ satisfy
\eqref{eq:rde} with $\xi\ge2$. Define
\begin{equation}\label{eq:mu-k}
 \mu=\E X,\qquad
 \widehat\mu_1=\frac{\E[X/(c+X)^2]}{\E[1/(c+X)^2]}.
\end{equation}
Let $\widehat\pi$ be the probability measure obtained by weighting the
law of $X$ with $(c+X)^{-2}$. A hat indicates this reweighting; write
\begin{equation}\label{eq:tilted-moments}
 \widehat\mu_j=\E_{\widehat\pi}X^j,\qquad j=1,2,3.
\end{equation}

\begin{lemma}[Two-child moment constraint]\label{lem:two-child}
The reweighted moments satisfy
\begin{align}
 \lambda(\widehat\mu_1-\widehat\mu_2)
 &\ge2\widehat\mu_1(1-2\widehat\mu_1+\widehat\mu_2),\notag\\
 \widehat\mu_1^2\le\widehat\mu_2&\le\overline\mu_2
 :=\frac{\widehat\mu_1(c-1+4\widehat\mu_1)}{c+1+2\widehat\mu_1}.
 \label{eq:moment-constraint}
\end{align}
Moreover,
\begin{equation}\label{eq:moment-feasible}
 \overline\mu_2-\widehat\mu_1^2
 =\frac{\widehat\mu_1(1-\widehat\mu_1)(c+2\widehat\mu_1-1)}
        {c+1+2\widehat\mu_1},
\end{equation}
and hence $\Delta:=c+2\widehat\mu_1-1\ge0$.
\end{lemma}

\begin{proof}
Let $\mathbb Q$ be the probability measure on $(\xi,X_1,X_2,\ldots)$ obtained by
weighting its original law with $(c+S)^{-2}$. Write
$S=X_1+S^{(1)}$, where $S^{(1)}=\sum_{i=2}^{\xi}X_i$ is the sum with the
first child omitted and is independent of $X_1$ under the original law.
The $X_1$-marginal under $\mathbb Q$, relative to $\widehat\pi$, has density proportional to
\begin{equation}\label{eq:tilt-ratio}
 R(x)=(c+x)^2\E(c+x+S^{(1)})^{-2}
     =\E\biggl[\biggl(\frac{c+x}{c+x+S^{(1)}}\biggr)^2\biggr].
\end{equation}
This is increasing in $x$. Consequently
\begin{equation}\label{eq:first-child}
 \E_{\mathbb Q}X_1=\frac{\E_{\widehat\pi}[XR(X)]}{\E_{\widehat\pi}R(X)}\ge \E_{\widehat\pi}X=\widehat\mu_1.
\end{equation}
For completeness, the covariance used here is nonnegative because, for
independent $U,V$ with law $\widehat\pi$,
$2\Cov_{\widehat\pi}(X,R(X))=\E[(U-V)(R(U)-R(V))]\ge0$.

Conditional exchangeability of the children now gives the decisive estimate
\begin{equation}\label{eq:two-children}
 \boxed{\quad
 \E_{\mathbb Q}S
 =\frac{\E[\xi X_1(c+S)^{-2}]}{\E[(c+S)^{-2}]}
 \ge2\E_{\mathbb Q}X_1\ge2\widehat\mu_1.
 \quad}
\end{equation}
The factor two in this inequality is the branching input of both proofs.
All quantities are integrable under this tilt, since
$S/(c+S)^2\le1/(4c)$.

By \eqref{eq:rde}, $S$ has the same law as $\lambda X/(1-X)$, and
$c+S$ has the same law as $(c+X)/(1-X)$. Thus
\begin{equation}\label{eq:QS}
 \E_{\mathbb Q}S=\lambda\frac{\E_{\widehat\pi}[X(1-X)]}{\E_{\widehat\pi}[(1-X)^2]}
     =\lambda\frac{\widehat\mu_1-\widehat\mu_2}{1-2\widehat\mu_1+\widehat\mu_2}.
\end{equation}
Combining \eqref{eq:two-children} and \eqref{eq:QS} proves
\eqref{eq:moment-constraint}. Jensen's inequality gives
$\widehat\mu_2\ge\widehat\mu_1^2$. The identity
\eqref{eq:moment-feasible} follows by subtraction, and its nonnegativity
implies $\Delta\ge0$.
\end{proof}

Returning from $\widehat\pi$ to the original law gives
\begin{equation}\label{eq:undo-tilt}
 \mu=
 \frac{\E_{\widehat\pi}[X(c+X)^2]}{\E_{\widehat\pi}[(c+X)^2]}
 =\frac{c^2\widehat\mu_1+2c\widehat\mu_2+\widehat\mu_3}
        {c^2+2c\widehat\mu_1+\widehat\mu_2}.
\end{equation}

\begin{remark}\label{rem:why-two}
The children need not remain independent under $\mathbb Q$. Independence is used
under the original measure in \eqref{eq:tilt-ratio}; conditional exchangeability
is preserved by the symmetric tilt. The lower bound
$\xi\ge2$ is needed in \eqref{eq:two-children}. Merely knowing a support bound
for $X$ does not supply either of these two pieces of information. If
single-child vertices are allowed, the same argument retains only
$\E_{\mathbb Q}S\ge \E_{\mathbb Q}X_1$ and does not give the factor two in Lemma~\ref{lem:two-child}.
\end{remark}

\subsection{Shared variation formulas}

The remaining algebraic preparation does not require a recursive fixed
point. Fix $\lambda>1$ and a probability law $\pi$ supported on $(0,1)$;
use the same inverse-square reweighting to define $\widehat\pi$ and
$\mu,\widehat\mu_j$. Abbreviate
\begin{equation}\label{eq:basepoint}
 c=\lambda-1,\quad N=N_\lambda(\pi),\quad Z=Z_\lambda(\pi),\quad v=\mathcal V_\lambda(\pi),
 \quad D=D_\lambda.
\end{equation}
Here $X,Y$ are independent with law $\pi$, and $0<v<1$. We use the affine function
$h_v$ to express the threshold $v/(1+v)$, and the weighted integral
operator $\cJ$ to collect the squared-denominator terms:
\begin{equation}\label{eq:J-h}
 h_v(y)=(1+v)y-v,
 \qquad
 (\cJ g)(x)=\E\frac{g(Y)}{D(x,Y)^2}.
\end{equation}
Here $g:[0,1]\to\mathbb R$ is bounded and measurable, and $x\in[0,1]$.

The derivative of one coordinate in the symmetric comparison is
\begin{equation}\label{eq:variation-kernel}
 \frac{\partial}{\partial x}
       \frac{xy-vq(x,y)}{D(x,y)}
 =\frac{y^2+c h_v(y)}{D(x,y)^2}.
\end{equation}
We write $\Psi=\cJ(y^2+c h_v)$ for the resulting averaged kernel.
In this single-coordinate identity, $v$ is held fixed. Each subsequent
derivative formula specifies the parameter and law at which this value
is evaluated.

Symmetry and a decreasing-weight covariance bound give
\begin{equation}\label{eq:v-mu-bound}
 \frac{v}{1+v}\le\frac\mu2,
 \qquad v\le\frac{\mu}{2-\mu}.
\end{equation}
Indeed, $N+Z=2\E[X/D(X,Y)]$, while for fixed $x$,
$y\mapsto D(x,y)^{-1}$ is decreasing. The covariance identity used in
\eqref{eq:first-child} therefore gives
$\E[Y/D(x,Y)]\le\mu\E[1/D(x,Y)]$. Multiplication by $x$ and averaging
yield $N\le\mu(N+Z)/2$ and hence \eqref{eq:v-mu-bound}.

We shall also compare inverse-square weights at different external
coordinates. Relative to $\widehat\pi$, the probability law obtained
by weighting $Y$ with $D(x,Y)^{-2}$ has density proportional to
$\{(c+y)/D(x,y)\}^2$. This ratio is increasing in $y$, since
\begin{equation}\label{eq:ratio-derivative}
 \frac{\partial}{\partial y}\frac{c+y}{D(x,y)}
 =\frac{\lambda x}{D(x,y)^2}\ge0.
\end{equation}
Thus every increasing test function has at least as large a mean under
the $D(x,Y)^{-2}$ weight as under $\widehat\pi$. These two elementary
reweighting observations will be used in both proofs.

\section{A first proof by fixed-point correction}\label{sec:correction}

We first use the common moment constraint to make the weighted mean
large enough for a positive variation kernel. We then add a residual
that vanishes at every recursive fixed point, so that both the law and
parameter variations have the required signs.

\subsection{A weighted-mean bound for the first proof}

The correction proof needs only the following consequence of the
common moment constraint.

\begin{lemma}[Inverse-square reweighting]\label{lem:tilt}
Under the assumptions of Lemma~\ref{lem:two-child},
$\mu<2\widehat\mu_1$.
\end{lemma}

\begin{proof}
If $\widehat\mu_1\ge1/2$, the conclusion follows from $\mu<1$. Suppose henceforth that
$0<\widehat\mu_1<1/2$. Since $\widehat\mu_3\le \widehat\mu_2$,
\begin{align}
 &2\widehat\mu_1(c^2+2c\widehat\mu_1+\widehat\mu_2)
   -(c^2\widehat\mu_1+2c\widehat\mu_2+\widehat\mu_3)\notag\\
 &\hspace{1em}\ge c^2\widehat\mu_1+4c\widehat\mu_1^2
                  +(2\widehat\mu_1-2c-1)\widehat\mu_2\notag\\
 &\hspace{1em}\ge c^2\widehat\mu_1+4c\widehat\mu_1^2
                  +(2\widehat\mu_1-2c-1)\overline\mu_2,
 \label{eq:gap-bound}
\end{align}
where the last step uses $2\widehat\mu_1-2c-1<0$ and \eqref{eq:moment-constraint}.
Recall $\Delta=c+2\widehat\mu_1-1\ge0$ from Lemma~\ref{lem:two-child}. A direct rearrangement gives
\begin{align}
 &c^2\widehat\mu_1+4c\widehat\mu_1^2
       +(2\widehat\mu_1-2c-1)\overline\mu_2\notag\\
 &\quad=\frac{\widehat\mu_1}{c+1+2\widehat\mu_1}\Bigl[
 \Delta^3+2\Delta^2+\{2+2\widehat\mu_1(1-2\widehat\mu_1)\}\Delta\notag\\
 &\hspace{10em}+2(1-\widehat\mu_1)(1-2\widehat\mu_1)\Bigr]>0.
 \label{eq:positive-polynomial}
\end{align}
Every term in brackets is nonnegative and the last is strictly positive.
Equations \eqref{eq:undo-tilt}, \eqref{eq:gap-bound} and
\eqref{eq:positive-polynomial} imply $\mu<2\widehat\mu_1$.
\end{proof}

\begin{remark}[Regular trees]\label{rem:regular}
For the deterministic $d$-ary tree, $\beta_\lambda=1-\lambda/d$ for
$0<\lambda<d$ and
$v(\lambda)=(d-\lambda)/(d+\lambda)$. The reweighted mean $\widehat\mu_1$ equals
$\E\beta_\lambda$ in this case, so the strict inequality of
Lemma~\ref{lem:tilt} also holds for a point-mass fixed point.
\end{remark}

\subsection{From the weighted mean to a positive kernel}\label{sec:kernel}

For the first proof, specialize the notation above to
$\pi=\pi_\lambda$ and $v=v(\lambda)$, with $1<\lambda<m$.

\begin{lemma}[Uniform positive margin]\label{lem:positive-kernel}
With $\widehat\mu_1$ as in Lemma~\ref{lem:tilt},
\begin{equation}\label{eq:threshold-k}
 \frac{v}{1+v}<\widehat\mu_1.
\end{equation}
In particular, writing $\eta=(1+v)\widehat\mu_1-v>0$,
\begin{equation}\label{eq:J-positive}
 (\cJ h_v)(x)\ge\frac{\eta}{\lambda^2}>0,
 \qquad 0\le x\le1.
\end{equation}
\end{lemma}

\begin{proof}
By \eqref{eq:v-mu-bound} and Lemma~\ref{lem:tilt},
\begin{equation}\label{eq:v-mu-k}
 \frac{v}{1+v}\le\frac\mu2<\widehat\mu_1.
\end{equation}
Let $\widehat\mu_1(x)=\E[Y/D(x,Y)^2]/\E[1/D(x,Y)^2]$, so that
$\widehat\mu_1(0)=\widehat\mu_1$. The increasing reweighting in
\eqref{eq:ratio-derivative}, applied to $y\mapsto y$, gives
$\widehat\mu_1(x)\ge\widehat\mu_1$. Since $D(x,y)\le\lambda$,
\[
 (\cJ h_v)(x)
 =\E[D(x,Y)^{-2}]\bigl((1+v)\widehat\mu_1(x)-v\bigr)
 \ge\frac{\eta}{\lambda^2}.
\]
\end{proof}



\subsection{Choose the correction from the variation kernel}

We construct the corrected comparison before computing its variations.
Throughout the remainder of this section, $\lambda$, $\pi=\pi_\lambda$
and $v=v(\lambda)$ remain fixed; only $\gamma$ and $\rho$ will vary. For a smooth
function $u$ on $[0,1]$, set
\begin{align}
 \res_u(\gamma,\rho)
 &=\E_\rho u(X)-\E_{\xi,\rho}u\bigl(f_\gamma(S_\rho)\bigr),
 &S_\rho&=\sum_{i=1}^{\xi}X_i,\label{eq:residual}\\
 \cF(\gamma,\rho)
 &=N_\gamma(\rho)-vZ_\gamma(\rho)-2\res_u(\gamma,\rho).
 \label{eq:corrected-functional}
\end{align}
Here $\gamma>1$, $\rho$ is a probability measure on $[0,1]$, and all
variables with law $\rho$ are independent of one another and of $\xi$.
The offspring law is unchanged. At every actual fixed point with
$1<\gamma<m$,
$\res_u(\gamma,\pi_\gamma)=0$, so
\begin{equation}\label{eq:on-curve}
 \boxed{\quad
 \cF(\gamma,\pi_\gamma)
 =Z_\gamma(\pi_\gamma)\{v(\gamma)-v(\lambda)\}.
 \quad}
\end{equation}
We seek a single choice of $u$ that makes $\cF$ increase with the law
and decrease with the explicit parameter near $(\lambda,\pi)$.

We use the single-coordinate derivative \eqref{eq:variation-kernel}.
To express the variation of the recursive term, define the linearization
operator $\cL$ on bounded functions $g$ by
\begin{equation}\label{eq:L}
 (\cL g)(x)=\E\biggl[
 \frac{\xi\lambda}{(\lambda+x+S^{(1)})^2}
 g\biggl(\frac{x+S^{(1)}}{\lambda+x+S^{(1)}}\biggr)
 \biggr],
\end{equation}
where the $X_i$ are independent with law $\pi$, independent of $\xi$,
and $S^{(1)}=\sum_{i=2}^{\xi}X_i$ again omits the first child.
Applied to $u'$, this operator records the contribution of all children when
one differentiates the recursive term in $\res_u$ along an interpolation
of the input variables. Conditional exchangeability gives the factor
$\xi$. Thus the variation of the residual is governed by $(I-\cL)u'$,
and that of $\cF$ by $\cJ(y^2+c h_v)-(I-\cL)u'$, where $I$ is the
identity operator. The exact transport identity is
proved in Lemma~\ref{lem:transport}.

Lemma~\ref{lem:positive-kernel} gives the sign of $\cJ h_v$. It therefore suffices
to choose $u$ so that $(I-\cL)u'=\cJ(y^2)$, cancelling the quadratic term.

Choose the derivative of the correction as
\begin{equation}\label{eq:r-u}
 \boxed{\quad u'(x)=(\cJ\mathrm{id})(x)=\E\frac{Y}{D(x,Y)^2}.\quad}
\end{equation}
Here $\mathrm{id}(y)=y$. The primitive normalized by $u(0)=0$ is
\begin{equation}\label{eq:u-explicit}
 u(x)=\int_0^x u'(s)\dd s
     =\E\frac{xY}{(c+Y)D(x,Y)}.
\end{equation}
Both $u$ and $u'$ are smooth and bounded on $[0,1]$.

\subsection{Verify the two signs}

\begin{proposition}[Cancellation of the distributional variation]\label{prop:cancel-law}
The correction \eqref{eq:r-u} satisfies
\begin{equation}\label{eq:Lr}
 \cL u'=\cJ\bigl(y(1-y)\bigr),\qquad
 (I-\cL)u'=\cJ(y^2).
\end{equation}
Consequently,
\begin{equation}\label{eq:law-margin}
 \boxed{\quad
 \bigl[\cJ(y^2+c h_v)-(I-\cL)u'\bigr](x)
 =c(\cJ h_v)(x)\ge\frac{c\eta}{\lambda^2}>0,
 \quad}
\end{equation}
uniformly for $x\in[0,1]$.
\end{proposition}

\begin{proof}
The elementary denominator identity
\begin{equation}\label{eq:denom-F}
 D(f_\lambda(s),y)
 =\frac{\lambda(c+s+y)}{\lambda+s}
\end{equation}
cancels the derivative of $f_\lambda$ against the squared kernel. Using an
additional independent copy $Y$ in \eqref{eq:L} gives
\begin{align}
 (\cL u')(x)
 &=\frac1\lambda\E\frac{\xi Y}{(c+x+S^{(1)}+Y)^2}\notag\\
 &=\frac1\lambda\E\frac{S}{(c+x+S)^2}\notag\\
 &=\E\frac{Y(1-Y)}{D(x,Y)^2}.\label{eq:Lr-proof}
\end{align}
For the middle equality, conditional on $\xi$, the $\xi$ variables
$Y,X_2,\ldots,X_\xi$ are exchangeable and their sum $Y+S^{(1)}$ has the
same law as $S$. The last equality uses \eqref{eq:rde}. Together with
$u'=\cJ\mathrm{id}$ and \eqref{eq:J-positive}, this gives
\eqref{eq:Lr} and \eqref{eq:law-margin}.
\end{proof}

\begin{proposition}[Cancellation of the parameter variation]\label{prop:cancel-parameter}
Define
\begin{equation}\label{eq:K}
 \kappa=\frac2\lambda\E[X(1-X)u'(X)]
        -\E\frac{v\,q(X,Y)-XY}{D(X,Y)^2}.
\end{equation}
Then
\begin{align}
 \lambda \kappa
 ={}&2\E\bigl[X(1-X)(\cJ h_v)(X)\bigr]\notag\\
 &+(1+v)\E\frac{XY(1-X)(1-Y)}{D(X,Y)^2}>0.
 \label{eq:K-positive}
\end{align}
\end{proposition}

\begin{proof}
By symmetry,
\begin{equation}\label{eq:K-expanded}
 \kappa=\E\frac{XY\{1+(2-X-Y)/\lambda\}-v\,q(X,Y)}{D(X,Y)^2}.
\end{equation}
Since $\lambda-D(x,y)=(1-x)(1-y)$ and
$\E[(XY-vq(X,Y))/D(X,Y)]=N-vZ=0$, multiplying
\eqref{eq:K-expanded} by $\lambda$ removes the term with a single power
of $D$ in the denominator. The remaining numerator satisfies the pointwise
identity
\begin{align}
 &(1-x)(1-y)\{xy-vq(x,y)\}+xy(2-x-y)\notag\\
 &\quad=x(1-x)h_v(y)+y(1-y)h_v(x)
             +(1+v)xy(1-x)(1-y).\label{eq:pointwise-K}
\end{align}
Divide by $D(x,y)^2$ and average to obtain \eqref{eq:K-positive}.
Both terms there are positive, by Lemma~\ref{lem:positive-kernel} and
$0<X,Y<1$ almost surely.
\end{proof}

\section{Closing the fixed-point correction proof}\label{sec:comparison}

The identities above give the required signs at the base point. To compare
two actual fixed points, the positive law-variation kernel must persist
along the entire interpolation between them. We establish this local
stability first, then change the law and the explicit parameter in turn.
The quantities $\lambda$, $\pi$, $v$ and $u$ remain fixed at the
base point \eqref{eq:basepoint}, with $\pi=\pi_\lambda$ and
$v=v(\lambda)$; in particular, $u$ is not recomputed
when $\gamma$ or $\rho$ changes.

\subsection{A transport formula and its local stability}

Write $\cL_{\gamma,\rho}$ for \eqref{eq:L} with parameter $\gamma$ and
independent child variables with law $\rho$. Define the variation kernel
\begin{equation}\label{eq:G-field}
 G_{\gamma,\rho}(x)
 =\int\frac{y^2+(\gamma-1)h_v(y)}{D_\gamma(x,y)^2}\,\rho(\mathrm dy)
   -u'(x)+(\cL_{\gamma,\rho}u')(x).
\end{equation}
At the base point, Proposition~\ref{prop:cancel-law} states that
\begin{equation}\label{eq:G-base}
 G_{\lambda,\pi}(x)=c(\cJ h_v)(x)
 \ge\frac{c\eta}{\lambda^2}>0,\qquad 0\le x\le1.
\end{equation}

\begin{lemma}[Transport and continuity]\label{lem:transport}
The map $(\gamma,\rho)\mapsto G_{\gamma,\rho}$ is continuous at
$(\lambda,\pi)$, with values in $C([0,1])$ equipped with the uniform norm and
with $\rho$ equipped with $W_1$.
Moreover, if $0\le U\le V\le1$ is a coupling, $X_s=(1-s)U+sV$, and
$\rho_s=\Law(X_s)$, then
\begin{equation}\label{eq:transport-formula}
 \frac{\mathrm d}{\mathrm ds}\cF(\gamma,\rho_s)
 =2\E\bigl[(V-U)G_{\gamma,\rho_s}(X_s)\bigr].
\end{equation}
Consequently, $\cF(\gamma,\cdot)$ is increasing along any such ordered
interpolation lying in a sufficiently small neighborhood of $\pi$, provided
$\gamma$ is sufficiently close to $\lambda$.
\end{lemma}

\begin{proof}
\smallskip\noindent
\underline{\textit{Step 1.}} Preserve positivity in a neighborhood.
Fix $1<a<\lambda<b<\infty$. The integral kernel in \eqref{eq:G-field} is
bounded and jointly continuous on $[a,b]\times[0,1]^2$.
For the recursive term, conditional on $\xi=n$ its integrand is bounded by
\begin{equation}\label{eq:nu-bound}
 \frac{n}{a}\|u'\|_\infty.
\end{equation}
For each fixed $n$ it is jointly continuous in the parameter, the external
coordinate $x$ and the $n-1$ child coordinates. Integrating against
$\rho^{\otimes(n-1)}$ therefore gives uniform convergence in $x$ when the
parameter and the law converge. One way to see uniformity is to use uniform
continuity of the integrand and then a finite net in $[0,1]$.
Truncating the sum at $n=M$ leaves a remainder bounded, uniformly in
$\gamma\in[a,b]$, $x$ and $\rho$, by
\begin{equation}\label{eq:nu-tail}
 \frac{\|u'\|_\infty}{a}\E[\xi\mathbf1_{\{\xi>M\}}]\longrightarrow0.
\end{equation}
This proves the asserted continuity. Notice that only $\E\xi<\infty$ is
used here. By \eqref{eq:G-base}, it follows that
$G_{\gamma,\rho}\ge c\eta/(2\lambda^2)$ uniformly in $x$ in a sufficiently small
neighborhood of $(\lambda,\pi)$.

\smallskip\noindent
\underline{\textit{Step 2.}} Compute the change along an interpolation.
Take independent copies of the coupled pair
$(U,V)$ for all child coordinates and for the two variables in
\eqref{eq:NZ}. Differentiation of $N_\gamma-vZ_\gamma$ gives twice the first
term of \eqref{eq:G-field}. Differentiation of $\E_{\rho_s}u(X)$ gives
$\E[(V-U)u'(X_s)]$. Finally, the chain rule for the RDE update produces
$\E[(V-U)(\cL_{\gamma,\rho_s}u')(X_s)]$ after conditional
exchangeability. Interchanging differentiation and expectation is justified
by \eqref{eq:nu-bound}, $0\le V-U\le1$ and $\E\xi<\infty$.
This proves \eqref{eq:transport-formula}.

The positivity from Step~1 makes \eqref{eq:transport-formula} nonnegative
on every ordered interpolation contained in that neighborhood.
\end{proof}

\subsection{The parameter derivative and the final comparison}

Freezing the distribution $\pi$ while varying the parameter gives
\begin{equation}\label{eq:parameter-residual}
 \left.\partial_\gamma\res_u(\gamma,\pi)\right|_{\gamma=\lambda}
 =\frac1\lambda\E\bigl[X(1-X)u'(X)\bigr].
\end{equation}
Indeed, at $\gamma=\lambda$,
$\partial_\gamma f_\gamma(s)=-s/(\lambda+s)^2
=-f_\lambda(s)(1-f_\lambda(s))/\lambda$.
The derivative is dominated since
$s/(\gamma+s)^2\le1/(4\gamma)$. From Proposition~\ref{prop:cancel-parameter},
\begin{equation}\label{eq:parameter-negative}
 \left.\partial_\gamma\cF(\gamma,\pi)\right|_{\gamma=\lambda}
 =-\kappa<0.
\end{equation}

\begin{proof}[First proof of \eqref{eq:main-dini}]
We first compare the two escape laws at the larger parameter, then use
the negative parameter derivative with the law frozen at $\pi$.

Let $\varepsilon>0$ be small enough that $\lambda+\varepsilon<m$.
By Lemma~\ref{lem:continuity}, on the same tree
\[
 U_\varepsilon:=\beta_{\lambda+\varepsilon}\le
 V:=\beta_\lambda,
 \qquad \E(V-U_\varepsilon)\longrightarrow0.
\]
The interpolation $X_s=(1-s)U_\varepsilon+sV$ satisfies
\begin{equation}\label{eq:interpolation-ball}
 W_1(\Law(X_s),\pi)
 \le\E|X_s-V|
 =(1-s)\E(V-U_\varepsilon)
 \le\E(V-U_\varepsilon).
\end{equation}
Thus, for sufficiently small $\varepsilon$, the entire interpolation is
inside the neighborhood in Lemma~\ref{lem:transport}. With
$\gamma=\lambda+\varepsilon$, that lemma gives
\begin{equation}\label{eq:law-comparison}
 \cF(\lambda+\varepsilon,\pi_{\lambda+\varepsilon})
 \le\cF(\lambda+\varepsilon,\pi).
\end{equation}
Since $\cF(\lambda,\pi)=0$, \eqref{eq:parameter-negative} yields
\[
 \cF(\lambda+\varepsilon,\pi)=-\kappa\varepsilon+o(\varepsilon).
\]
Using \eqref{eq:on-curve}, we conclude that
\begin{equation}\label{eq:finite-comparison}
 Z_{\lambda+\varepsilon}(\pi_{\lambda+\varepsilon})
 \{v(\lambda+\varepsilon)-v(\lambda)\}
 \le-\kappa\varepsilon+o(\varepsilon).
\end{equation}
The normalizing factor converges to $Z>0$. Therefore
\begin{equation}\label{eq:explicit-dini}
 \boxed{\quad \Dplus v(\lambda)\le-\frac{\kappa}{Z}<0.\quad}
\end{equation}
This proves \eqref{eq:main-dini}, with $\delta_\lambda=\kappa/Z$.

\end{proof}

\section{A second proof by stochastic order and root response}\label{sec:root-proof}

We give a second proof of \eqref{eq:main-dini} by comparing the actual
escape law with the law obtained when the parameter changes only at the
root. This comparison uses the original ratio
$\mathcal V_\gamma(\rho)=N_\gamma(\rho)/Z_\gamma(\rho)$.
The same two-child moment constraint now supplies a stronger inequality:
it makes $\mathcal V$ locally increasing under stochastic order and makes
the root response strictly negative.

\subsection{A stronger consequence of the tilted-moment constraint}

\begin{proposition}[A strict tilted-moment inequality]\label{prop:root-Gamma}
Fix $\lambda>1$, put $c=\lambda-1$, and let $X\in(0,1)$ satisfy
\eqref{eq:rde} with $\xi\ge2$. With $\mu$, $\widehat\pi$ and
$\widehat\mu_j$ as in Lemma~\ref{lem:two-child},
\begin{equation}\label{eq:root-Gamma}
 \Gamma:=\lambda(2\widehat\mu_1-\mu)
       -(1-\mu)(\widehat\mu_1-\widehat\mu_2)>0.
\end{equation}
\end{proposition}

\begin{proof}
If $\widehat\mu_1\ge1/2$, then
\[
 \Gamma=\lambda(2\widehat\mu_1-1)
       +(1-\mu)(\lambda-\widehat\mu_1+\widehat\mu_2)>0.
\]
The first term is nonnegative, and the second is strictly positive since
$\mu<1$ and $\lambda>1$.

Suppose $0<\widehat\mu_1<1/2$. By \eqref{eq:undo-tilt} and
$\widehat\mu_3\le\widehat\mu_2$,
\begin{equation}\label{eq:root-mu-upper}
 \mu\le\mu_+(\widehat\mu_2),\qquad
 \mu_+(s):=
 \frac{c^2\widehat\mu_1+(2c+1)s}
      {c^2+2c\widehat\mu_1+s}.
\end{equation}
The expression defining $\Gamma$ decreases strictly with $\mu$, whose
coefficient is $-(\lambda-\widehat\mu_1+\widehat\mu_2)<0$.
It therefore suffices to show that
\begin{equation}\label{eq:root-F}
 \begin{aligned}
 F(s):={}&(c^2+2c\widehat\mu_1+s)\Bigl[
  (c+1)(2\widehat\mu_1-\mu_+(s))\Bigr.\\
 &\hspace{8em}\Bigl.
  -(1-\mu_+(s))(\widehat\mu_1-s)\Bigr]
 \end{aligned}
\end{equation}
is positive at $s=\widehat\mu_2$.
After clearing the denominator in $\mu_+(s)$, $F$ is a quadratic
polynomial with leading coefficient $-2c$. It is concave. The shared
constraint \eqref{eq:moment-constraint} and Jensen's inequality give
$\widehat\mu_1^2\le\widehat\mu_2\le\overline\mu_2$, so positivity at
the two endpoints is sufficient.

Recall $\Delta=c+2\widehat\mu_1-1\ge0$ from
\eqref{eq:moment-feasible}. Direct expansion gives
\begin{equation}\label{eq:root-F-endpoints}
 F(\widehat\mu_1^2)=\widehat\mu_1P_3(\Delta,\widehat\mu_1),\qquad
 F(\overline\mu_2)
 =\frac{\widehat\mu_1P_5(\Delta,\widehat\mu_1)}{(\Delta+2)^2},
\end{equation}
where
\begin{align}
 P_3(\Delta,t)
 ={}&\Delta^3+(1-t)(t+3)\Delta^2\notag\\
 &+(1-t)(3-2t-2t^2)\Delta+(1-t)^2(1-2t),
 \label{eq:root-P3}\\
 P_5(\Delta,t)
 ={}&\Delta^5+2(3-t)\Delta^4+2(6-t-3t^2)\Delta^3\notag\\
 &+2(1-t)(5+12t-8t^2)\Delta^2\notag\\
 &+2(1-t)(3+7t-16t^2+4t^3)\Delta\notag\\
 &+4(1-t)^2(1-2t).
 \label{eq:root-P5}
\end{align}
For $0<t<1/2$, all these coefficients are positive. Indeed,
$3-2t-2t^2>3/2$, $6-t-3t^2>19/4$ and $5+12t-8t^2>0$, while
\[
 3+7t-16t^2+4t^3=3+t(1-2t)(7-2t)>0.
\]
Both constant terms are strictly positive, so the two polynomials are
positive for $\Delta\ge0$. Equation~\eqref{eq:root-F-endpoints} and
concavity give $F(\widehat\mu_2)>0$. Together with
\eqref{eq:root-mu-upper}, this proves \eqref{eq:root-Gamma}.
\end{proof}

\subsection{The distributional condition behind the comparison}

For the next two comparisons, fix $\lambda>1$ and an arbitrary probability
law $\pi$ supported on $(0,1)$, with the notation in
\eqref{eq:basepoint}. In particular, $X,Y$ are independent with law $\pi$
and $v=\mathcal V_\lambda(\pi)$; the moments $\mu$ and
$\widehat\mu_j$ are computed using this law. No recursive equation is
assumed in these distributional comparisons.

The function needed for the root response is
\begin{equation}\label{eq:root-h-tilde}
 \widetilde h_v(y)=h_v(y)-\frac{1-v}{2\lambda}y(1-y).
\end{equation}
The extra term in $\widetilde h_v$ will arise from the cancellation of
the root transport with the explicit parameter change.

\begin{lemma}[Positive kernels from the scalar condition]\label{lem:root-kernel}
If $\Gamma$ defined by \eqref{eq:root-Gamma} is positive, then, for every
$x\in[0,1]$,
\begin{align}
 (\cJ\widetilde h_v)(x)
 &\ge\frac{\Gamma}{\lambda^3(2-\mu)}>0,
 \label{eq:root-h-positive}\\
 \Psi(x)&\ge\frac{c\Gamma}{\lambda^3(2-\mu)}>0.
 \label{eq:root-Psi-positive}
\end{align}
\end{lemma}

\begin{proof}
The shared bound \eqref{eq:v-mu-bound} gives $v\le\mu/(2-\mu)$.
For $0\le y\le1$,
\[
 \partial_v\widetilde h_v(y)
 =-(1-y)\left(1-\frac{y}{2\lambda}\right)\le0,
 \qquad
 \widetilde h_v'(y)
 =1+v-\frac{1-v}{2\lambda}(1-2y)>0.
\]
Consequently,
\begin{equation}\label{eq:root-tilted-average}
 \E_{\widehat\pi}\widetilde h_v(X)
 \ge\E_{\widehat\pi}\widetilde h_{\mu/(2-\mu)}(X)
 =\frac{\Gamma}{\lambda(2-\mu)}.
\end{equation}

For each $x$, the law tilted by $D(x,Y)^{-2}$, relative to
$\widehat\pi$, has density proportional to
$\{(c+y)/D(x,y)\}^2$. This density is increasing by
\eqref{eq:ratio-derivative}. The covariance inequality therefore shows
that the mean of the increasing function $\widetilde h_v$ under this
tilt is at least the mean in
\eqref{eq:root-tilted-average}. Since $D(x,y)\le\lambda$, the total
weight is at least $\lambda^{-2}$, proving
\eqref{eq:root-h-positive}. Finally,
\[
 y^2+c h_v(y)
 =c\widetilde h_v(y)+y^2+\frac{c(1-v)}{2\lambda}y(1-y).
\]
The two last terms are nonnegative, giving
\eqref{eq:root-Psi-positive}.
\end{proof}

\begin{lemma}[A: local stochastic-order comparison]\label{lem:root-A}
If $\Gamma>0$, there exist $\eta,r>0$ such that, whenever $\gamma>1$
and probability laws $\rho_0,\rho_1$ on $[0,1]$ satisfy
$|\gamma-\lambda|<\eta$ and $W_1(\rho_i,\pi)<r$, $i=0,1$,
\begin{equation}\label{eq:root-A}
 \rho_0\le_{\mathrm{st}}\rho_1
 \quad\Longrightarrow\quad
 \mathcal V_\gamma(\rho_0)\le\mathcal V_\gamma(\rho_1).
\end{equation}
The inequality is strict when $\rho_0\ne\rho_1$.
\end{lemma}

\begin{proof}
For $\gamma>1$ and $\rho\ne\delta_0$, put
$v_{\gamma,\rho}=\mathcal V_\gamma(\rho)$ and
\begin{equation}\label{eq:root-Psi-general}
 \Psi_{\gamma,\rho}(x)
 =\int\frac{y^2+(\gamma-1)((1+v_{\gamma,\rho})y-v_{\gamma,\rho})}
              {D_\gamma(x,y)^2}\,\rho(\mathrm dy).
\end{equation}
This kernel depends continuously on $(\gamma,\rho)$ at $(\lambda,\pi)$
in the uniform norm on $[0,1]$, with laws equipped with $W_1$.
Indeed, on a compact parameter interval in $(1,\infty)$, the kernels
are jointly continuous with denominators bounded away from zero.
Their product integrals are continuous in $\rho$, and
$Z_\lambda(\pi)>0$ keeps the ratio $v_{\gamma,\rho}$ well defined nearby.
Uniform continuity in $x$ gives the uniform-norm assertion.
Thus \eqref{eq:root-Psi-positive} persists: for some
$\eta,r,\delta>0$,
\begin{equation}\label{eq:root-uniform-neighborhood}
 \Psi_{\gamma,\rho}(x)\ge\delta
 \quad\hbox{for all }x\in[0,1]
\end{equation}
when $|\gamma-\lambda|<\eta$ and $W_1(\rho,\pi)<r$.
We take $r$ small enough to exclude $\delta_0$.

Couple the ordered laws by their quantiles: if $U$ is uniform on $(0,1)$,
set $X_i=F_{\rho_i}^{-1}(U)$, where $F_{\rho_i}$ is the
distribution function of $\rho_i$. Then $X_0\le X_1$. For
$X_s=(1-s)X_0+sX_1$ and $\rho_s=\Law(X_s)$,
\begin{equation}\label{eq:root-quantile-ball}
 W_1(\rho_s,\pi)
 \le(1-s)W_1(\rho_0,\pi)+sW_1(\rho_1,\pi)<r.
\end{equation}
Using an independent copy of this coupled interpolation for the second
coordinate, the single-coordinate calculation
\eqref{eq:variation-kernel} and the quotient rule give
\begin{equation}\label{eq:root-transport}
 \frac{\mathrm d}{\mathrm ds}\mathcal V_\gamma(\rho_s)
 =\frac{2}{Z_\gamma(\rho_s)}
   \E[(X_1-X_0)\Psi_{\gamma,\rho_s}(X_s)].
\end{equation}
In applying the single-coordinate calculation at each $s$, the scalar
$v=v_{\gamma,\rho_s}$ is held fixed. All differentiated kernels are
bounded, which justifies differentiation under the integrals.
Equations~\eqref{eq:root-uniform-neighborhood} and
\eqref{eq:root-quantile-ball} make the derivative nonnegative. If the
two laws differ, $\E(X_1-X_0)>0$, so the derivative is strictly positive.
Integration over $s\in[0,1]$ proves Lemma A.
\end{proof}

\begin{samepage}
\begin{lemma}[B: strictly negative root response]\label{lem:root-B}
Suppose $\Gamma>0$. Define
\begin{equation}\label{eq:root-map}
 T_{\lambda,\varepsilon}(x)
 =\frac{\lambda x}{\lambda+\varepsilon(1-x)},\qquad
 H_\lambda(\varepsilon)
 =\mathcal V_{\lambda+\varepsilon}
       ((T_{\lambda,\varepsilon})_\#\pi).
\end{equation}
Then $H_\lambda$ is differentiable at $0$, and
\begin{equation}\label{eq:root-kappa-definition}
 \begin{aligned}
 \kappa_{\mathrm{root}}
 &:=\frac2\lambda\E[X(1-X)\Psi(X)]
       -\E\frac{v\,q(X,Y)-XY}{D(X,Y)^2}\\
 &=2\E[X(1-X)(\cJ\widetilde h_v)(X)]>0.
 \end{aligned}
\end{equation}
More precisely,
\begin{equation}\label{eq:root-negative}
 H_\lambda'(0)=-\frac{\kappa_{\mathrm{root}}}{Z}
 \le-\frac{2\Gamma\E[X(1-X)]}{Z\lambda^3(2-\mu)}<0.
\end{equation}
\end{lemma}
\end{samepage}

\begin{proof}
For $\varepsilon$ in a small interval about $0$, the parameter is greater
than $1$, the map $T_{\lambda,\varepsilon}$ sends $[0,1]$ into itself,
and its derivatives and the kernel derivatives are bounded.
At $\varepsilon=0$ its velocity is
\[
 \left.\partial_\varepsilon T_{\lambda,\varepsilon}(x)
       \right|_{\varepsilon=0}=-\frac{x(1-x)}{\lambda}.
\]
The explicit parameter derivative of $N-vZ$, with $v$ and $\pi$ held
fixed, is $\E[(v q(X,Y)-XY)/D(X,Y)^2]$.
The two transported coordinates contribute
$-2\E[X(1-X)\Psi(X)]/\lambda$ by
\eqref{eq:variation-kernel}. The quotient rule therefore gives
$H_\lambda'(0)=-\kappa_{\mathrm{root}}/Z$.

To see the sign, expand the following identity:
\begin{align}
 &x(1-x)\{y^2+c h_v(y)\}+y(1-y)\{x^2+c h_v(x)\}
       +\lambda\{xy-vq(x,y)\}\notag\\
 &\qquad=\lambda\{x(1-x)\widetilde h_v(y)
                   +y(1-y)\widetilde h_v(x)\}
       +D(x,y)\{xy-vq(x,y)\}.
 \label{eq:root-pointwise-cancellation}
\end{align}
Divide by $\lambda D(x,y)^2$ and integrate. The final term has mean
zero since $N-vZ=0$, and symmetry gives the second equality in
\eqref{eq:root-kappa-definition}. Lemma~\ref{lem:root-kernel} now gives
\eqref{eq:root-negative}; it is strict because
$\E[X(1-X)]>0$.
\end{proof}

These comparisons use only $\Gamma>0$ for the law $\pi$.
Proposition~\ref{prop:root-Gamma} supplies this condition at the escape
law; the local stochastic-order comparison and the root-response bound
themselves have no fixed-point assumption.

\subsection{Comparing the root response with the actual escape law}

We return to the tree model and fix $1<\lambda<m$. In
\eqref{eq:basepoint}, take $\pi=\pi_\lambda$ and $v=v(\lambda)$. Increasing the parameter decreases the escape
probability in every child subtree. Freezing these child probabilities
therefore gives an upper bound on the escape probability at the root.

\begin{lemma}[Freezing the descendants]\label{lem:root-freezing}
For $\varepsilon>0$ with $\lambda+\varepsilon<m$, the same-tree coupling
satisfies
\begin{equation}\label{eq:root-freezing}
 \beta_{\lambda+\varepsilon}
 \le T_{\lambda,\varepsilon}(\beta_\lambda),\qquad
 \pi_{\lambda+\varepsilon}
 \le_{\mathrm{st}}(T_{\lambda,\varepsilon})_\#\pi_\lambda.
\end{equation}
\end{lemma}

\begin{proof}
Let $S_\gamma$ be the sum of the child escape probabilities at parameter
$\gamma$. Lemma~\ref{lem:continuity} gives
$S_{\lambda+\varepsilon}\le S_\lambda$, so
\[
 \beta_{\lambda+\varepsilon}
 =\frac{S_{\lambda+\varepsilon}}
        {\lambda+\varepsilon+S_{\lambda+\varepsilon}}
 \le\frac{S_\lambda}{\lambda+\varepsilon+S_\lambda}.
\]
The base-point recursion gives
$S_\lambda=\lambda\beta_\lambda/(1-\beta_\lambda)$.
Substitution identifies the last expression with
$T_{\lambda,\varepsilon}(\beta_\lambda)$ and proves both claims.
\end{proof}

\begin{proof}[Second proof of \eqref{eq:main-dini}]
Proposition~\ref{prop:root-Gamma} gives $\Gamma>0$ at
$\pi=\pi_\lambda$, so Lemmas~\ref{lem:root-A}
and~\ref{lem:root-B} apply there.
By Lemma~\ref{lem:continuity},
$W_1(\pi_{\lambda+\varepsilon},\pi)\to0$.
The other law in \eqref{eq:root-freezing} is also close to $\pi$, since
\[
 0\le x-T_{\lambda,\varepsilon}(x)
 =\frac{\varepsilon x(1-x)}{\lambda+\varepsilon(1-x)}
 \le\frac{\varepsilon}{4\lambda}.
\]
For sufficiently small $\varepsilon>0$, both laws and the parameter
$\lambda+\varepsilon$ thus lie in the joint neighborhood in
Lemma~\ref{lem:root-A}. Applying Lemma A at the
perturbed parameter gives
\begin{equation}\label{eq:root-closing-comparison}
 \begin{aligned}
 v(\lambda+\varepsilon)
 &=\mathcal V_{\lambda+\varepsilon}(\pi_{\lambda+\varepsilon})\\
 &\le\mathcal V_{\lambda+\varepsilon}
       ((T_{\lambda,\varepsilon})_\#\pi)
 =H_\lambda(\varepsilon).
 \end{aligned}
\end{equation}
Because $H_\lambda(0)=v(\lambda)$,
Lemma~\ref{lem:root-B} yields
\begin{equation}\label{eq:root-dini}
 \Dplus v(\lambda)
 \le-\frac{\kappa_{\mathrm{root}}}{Z}
 \le-\frac{2\Gamma\E[\beta_\lambda(1-\beta_\lambda)]}
             {Z\lambda^3(2-\E\beta_\lambda)}<0.
\end{equation}
Only bounded kernels at a fixed law were differentiated; the change in
the actual fixed-point law was handled by the finite comparison
\eqref{eq:root-closing-comparison}.
\end{proof}

\begin{remark}[The two descent bounds]\label{rem:root-two-bounds}
The correction argument gives \eqref{eq:explicit-dini}, while the root
comparison gives \eqref{eq:root-dini}. Both use the same normalizing
factor $Z$, so together they give
$\Dplus v(\lambda)\le-\max\{\kappa,\kappa_{\mathrm{root}}\}/Z$.
The first bound uses the residual identity along the fixed-point curve;
the second uses a stronger tilted-moment estimate to make the root
response alone strictly negative.
\end{remark}

\clearpage
\section{Completing the speed theorem}\label{sec:completion}

Either proof gives \eqref{eq:main-dini} on $(1,m)$. We now make the
same continuity and low-parameter arguments for both routes.

\subsection{From a negative Dini derivative to strict decrease}

A continuous function with strictly negative upper right Dini derivative at
every point of an interval is strictly decreasing there. Indeed, if
$f(b)\ge f(a)$ for $a<b$, a minimum of $f$ on $[a,b]$ is attained at some
$x<b$: if $b$ is a minimizer, then $a$ is also one. At such an $x$ all
sufficiently small right difference quotients are nonnegative, a
contradiction. By Lemma~\ref{lem:continuity}, this argument applies to $v$ on
$(1,m)$ using \eqref{eq:main-dini}.

\subsection{The low-parameter interval and classical derivatives}

Let $m_1=\min\{k:p_k>0\}\ge2$. Song--Wang--Xiang
\cite[Theorem~1.2 and Lemma~2.5]{SWX} establish $v'(\lambda)<0$ for
\begin{equation}\label{eq:low-range}
 0<\lambda\le a_{m_1},\qquad
 a_{m_1}:=\frac{m_1}{1+\sqrt{1-1/m_1}}>1,
\end{equation}
and strict decrease on $[0,a_{m_1}]$. Their derivative statement and its
justification are given immediately before, and in, Lemma~2.5 of the cited
preprint. This interval overlaps $(1,m)$, so its strict monotonicity and
\eqref{eq:main-dini} complete Theorem~\ref{thm:main} on $[0,m)$.

Under the exponential-moment assumption of Corollary~\ref{cor:derivative},
\cite[Theorem~1.1]{BT} gives differentiability on $(0,m)$: its lower endpoint
is zero in the leafless case. At every $\lambda>1$ the derivative equals the
Dini derivative and is negative by \eqref{eq:main-dini}. The remaining
parameters are covered by \eqref{eq:low-range}. This proves
Corollary~\ref{cor:derivative}.

\begin{thebibliography}{99}

\bibitem[A\"id14]{Aidekon}
\'{E}.~A\"id\'ekon.
\emph{Speed of the biased random walk on a Galton--Watson tree}.
\emph{Probab. Theory Related Fields} \textbf{159} (2014), 597--617.
\href{https://doi.org/10.1007/s00440-013-0515-y}{doi:10.1007/s00440-013-0515-y}.
The cited theorem numbering follows
\href{https://arxiv.org/abs/1111.4313v4}{arXiv:1111.4313v4}.

\bibitem[BT20]{BT}
A.~Bowditch and Y.~Tokushige.
\emph{Differentiability of the speed of biased random walks on
Galton--Watson trees}.
\emph{ALEA Lat. Am. J. Probab. Math. Stat.} \textbf{17} (2020), 609--642.
\href{https://doi.org/10.30757/ALEA.v17-24}{doi:10.30757/ALEA.v17-24}.
The cited theorem numbering follows
\href{https://arxiv.org/abs/1906.07913v2}{arXiv:1906.07913v2}.

\bibitem[LPP97]{LPP}
R.~Lyons, R.~Pemantle and Y.~Peres.
\emph{Unsolved problems concerning random walks on trees}.
In \emph{Classical and Modern Branching Processes},
IMA Vol. Math. Appl. \textbf{84}, 223--237.
Springer, New York, 1997.
\href{https://rdlyons.pages.iu.edu/pdf/opentree7.pdf}{Corrected author version}.

\bibitem[MK26]{MK}
M.~Mandarapu and S.~Kunkunuru.
\emph{A computer-assisted proof of speed monotonicity for the biased
random walk on a Galton--Watson tree beyond the known range}.
\href{https://arxiv.org/abs/2609.29894v1}{arXiv:2609.29894v1} (2026).

\bibitem[SWX25]{SWX}
H.~Song, L.~Wang and K.~Xiang.
\emph{The speed of a biased walk on a Galton--Watson tree without leaves
is monotonic for low values of bias}.
\emph{J. Appl. Probab.} \textbf{62} (2025), no.~3, 1044--1052.
\href{https://doi.org/10.1017/jpr.2024.113}{doi:10.1017/jpr.2024.113}.
The cited theorem and lemma numbering follows
\href{https://arxiv.org/abs/1610.08151v1}{arXiv:1610.08151v1}.

\end{thebibliography}

\end{document}
